\section{Implementering}

\emph{Implementeringen} realiserer den centrale model i TræningsMester. Planlagt træning, faktisk logging, completion, progression og sikker dataadgang skal hænge sammen. Fokus ligger på de flows, der bærer problemformuleringen.

Den praktiske test for implementeringen er, om brugerens træningsforløb kan fortsætte efter ændringer. En app kan have mange skærme og stadig fejle, hvis en afbrudt session, en ændret øvelse eller en manglende detaljelog gør progressionen uklar. Derfor betyder sammenhængen mellem events, state og backendregler mest.

\subsection{Session, repositories og programstruktur}

Appen etablerer session, profil og dataadgang, før brugerbundne træningsdata vises. SwiftUI gør brugerfladen state-drevet, men adgang og datakontrakter ligger i AppState, repositories og backendregler \cite{apple_swiftui_docs}. Supabase RLS og \emph{API-sikring} betyder, at klienten ikke alene afgør, hvad brugeren kan læse og skrive \cite{supabase_rls_docs,supabase_securing_api_docs}.

Programstrukturen giver den planlagte ramme. Programmer, workouts og øvelser beskriver, hvad brugeren forventes at gennemføre. Figur~\ref{fig:tm-data-model-er-left} viser denne del af datamodellen. Den planlagte struktur skal suppleres af en model for faktisk træning, fordi brugeren kan ændre øvelser, belastning eller logging undervejs.

\emph{Repositories} gør datatilgangen mere kontrolleret. I stedet for at hver visning skriver direkte mod backend, samles læsning, skrivning og mapping i et lag, der kan testes og ændres. Det passer til sikkerhedsargumentet, fordi klienten stadig er en klient, mens adgangsregler og privilegerede handlinger ligger uden for visningslaget.

\subsection{Tracker, completion og cycle runtime}

\emph{Tracker- og completion-flowet} er implementeringens svar på afvigelsesproblemet. Figur~\ref{fig:tm-ios-tracker-active} viser aktiv tracker med planlagte sæt, mens Figur~\ref{fig:tm-tracker-completion-summary} viser, at completion kan registrere gennemført træning uden fuld trackerlog. Det bevarer historik og næste handling, selv om datadetaljen er lavere.

\emph{Selvmonitorering} og feedback er centrale adfærdsændringsteknikker \cite{michie2013_bct_taxonomy_v1}. Feedbacken skal samtidig være let nok til brug under træning. Completion fungerer derfor som et bevidst lav-friktionsflow.

\emph{Cycle runtime} forbinder completion med næste workout. Appen kan bruge runtime-status til at afgøre, hvad der er relevant nu, i stedet for kun at gå videre i en fast planrækkefølge. Autoregulering gør denne fleksibilitet fagligt meningsfuld \cite{greig2020_autoregulation_resistance_training,larsen2021_autoregulation_systematic_review}. TræningsMester strukturerer næste handling ud fra faktisk gennemførelse.

Det er her implementeringen adskiller sig fra en simpel checklistelogik. Completion fungerer som et event, der kan påvirke cyklus, historik og Home-visning. Trackerlog giver et mere detaljeret datagrundlag for feedback. Begge dele skal kunne eksistere uden at modsige hinanden. Dermed kan appen reagere på brugerens faktiske forløb frem for kun at vise en liste over øvelser.

\subsection{Støtteflader og privilegerede flows}

\emph{watchOS} og \emph{Live Activity} implementeres som støtteflader til samme træningsmodel. De kan give hurtig status og adgang under træning, og de skal pege på samme session og completion som hovedappen. Apple beskriver Live Activities som tidsfølsom status og Watch Connectivity som koordinering mellem iPhone og Apple Watch \cite{apple_activitykit_live_data_docs,apple_watchconnectivity_docs}. Derfor genbruger støttefladerne samme sessions- og datakontrakt som hovedappen.

Edge Functions afgrænser flows, der kræver server-side logik. Det gælder blandt andet admin, import, AI, billing og træner-/klientrelationer. OWASP MASVS og Supabase API-sikring understøtter, at privilegerede nøgler og bredere databaseoperationer holdes server-side \cite{owasp2024_masvs,supabase_securing_api_docs}. Det centrale implementeringsbidrag er stadig sammenhængen mellem plan, tracker, completion, cyklusprogression og lav-friktionsflader.

Implementeringen viser dermed en prioritering. Kerneflowet skal være robust, før støttefunktioner bliver afgørende. Admin, import og betaling kan udvides senere, når de understøtter den basale model for brugerens træning. Derfor er test og verifikation især rettet mod build, domænelogik, sikkerhedsgrænser og kendte gaps.

\tmdelkonklusion{Implementeringens pointe}{Kerneideen er implementeret som events og fælles state. En completion kan få betydning for historik, cyklus og Home, mens en trackerlog kan give mere detaljeret feedback uden at være eneste vej gennem systemet.}
